\vspace*{\fill}

% Image 1 — AR 1.375 — full width
\begin{center}\bwimage[width=0.85\linewidth]{acte_06/chapter_14/_0000_origines-de-l-homme-moderne_2.jpg}\end{center}\vspace{-0.3em}
{\figcaptionname\itshape Jebel Irhoud, à une centaine de kilomètres à l'ouest de Marrakech
— le site marocain daté à \~300\,000 ans qui livre le plus ancien témoignage connu
d'\textit{Homo sapiens} sur le continent africain.}

{\centering\rule{0.4\textwidth}{0.4pt}\par}
\medskip  

\noindent\bwimage[width=0.95\linewidth]{acte_06/chapter_14/_0000_origines-de-l-homme-moderne.jpg}\\[0.3em]
% Image 2 — AR 1.587 — full width
{\figcaptionname\itshape Lieux et dates des principaux sites fossiles d'hominidés en
Afrique, illustrant la coexistence de plusieurs espèces — H.\,sapiens, H.\,naledi et H.\,heidelbergensis

\vspace*{\fill}
